\section*{Entire original Section 8: An easy polynomial with many vertices (pp.9:19--21)}
Here we construct a bivariate polynomial with a monotone arithmetic circuit of linear size, but whose Newton polytope has exponentially many vertices. This proves Theorem 30.
We use the following notation. Given $(a,b)\in\mathbb R^2$,
\[(a,b)\cdot P:=\{(ax,by):(x,y)\in P\}.\]
Given $a\in\mathbb R$,
\[aP:=(a,a)\cdot P.\]
\begin{sourceresult}{Observation 52}
For a bivariate polynomial $f(x,y)$,
\[\Newt(f(x^a,y^b))=(a,b)\Newt(f(x,y))\quad\text{and}\quad\Newt(f^a)=a\Newt(f).\]
\end{sourceresult}
The building block of the polynomial are the following two polytopes. Let $P_n$ be the polytope with vertices $\{(k,k^2):0\le k\le n-1\}$. Let $Q_n$ be the polytope with vertices $\{(k,k^2+k):0\le k\le n-1\}$. These polytopes can be constructed inductively as follows.
\begin{sourceresult}{Lemma 53}
For every $n\ge1$,
\begin{align}
P_{2n}&=(2,4)\cdot P_n\sqcup((1,1)+(2,4)\cdot Q_n)),\tag{8}\\
Q_{2n}&=(1,2)\cdot(P_n+Q_n)\sqcup\{(2n-1,2n(2n-1))\}.\tag{9}
\end{align}
\end{sourceresult}
\begin{proof}
Part (8). Let $0\le k\le2n-1$. If $k=2r$ is even then $r\le n-1$ and
\[(k,k^2)=(2,4)(r,r^2)\]
with $(r,r^2)$ a vertex of $P_n$. If $k=2r+1$ is odd then $r\le n-1$ and
\[(k,k^2)=(2r+1,4r^2+4r+1)=(1,1)+(2,4)\cdot(r,r^2+r),\]
where $(r,r^2+r)$ is a vertex of $Q_n$. This shows the containment $\subseteq$ in (8). The other direction holds since $P_n\sqcup Q_n$ can have at most $2n$ vertices.

Part (9). We first describe the vertices of $(1,2)(P_n+Q_n)$. We claim that
\begin{equation}
\vertset((1,2)(P_n+Q_n))=\{v_0,v_1,\ldots,v_{2n-2},u\},\tag{10}
\end{equation}
where $v_k:=(k,k^2+k)$, $u:=(n-1,2n(n-1))$.
Given $0\le k\le2n-2$, let us show that $v_k$ is a vertex of $(1,2)(P_n+Q_n)$. If $k=2r$ is even, we have $r\le n-1$ and
\[(k,k^2+k)=(2r,4r^2+2r)=(1,2)(r,r^2)+(1,2)(r,r^2+r).\]
If $k=2r+1$ is odd, we have $r\le n-2$ and
\[(k,k^2+k)=(2r+1,4r^2+6r+2)=(1,2)(r+1,(r+1)^2)+(1,2)(r,r^2+r).\]
This means that $v_k\in(1,2)(P_n+Q_n)$. Now, every $(z_1,z_2)\in(1,2)(P_n+Q_n)$ satisfies $z_2\ge z_1^2+z_1$, because
\[2r_1^2+2(r_2^2+r_2)-(r_1+r_2)^2-(r_1+r_2)=(r_1-r_2)^2-(r_1-r_2)\ge0.\]
Since $v_k$ lies on the curve $z_2=z_1^2+z_1$, and the curve is strictly convex, $v_k$ cannot be convex combination of other points in $(1,2)(P_n+Q_n)$. So, $v_k$ is indeed a vertex. To show that $u$ is a vertex, note that both $(1,2)P_n$ and $(1,2)Q_n$ are contained in the halfplane $\{(z_1,z_2)\in\mathbb R^2:z_2\le2nz_1\}$. On the boundary $z_2=2nz_1$, $(1,2)Q_n$ has vertices $(0,0)$ and $u$, and $(1,2)P_n$ only the vertex $(0,0)$. This implies $u$ is a vertex of $(1,2)(P_n+Q_n)$. This proves the containment $\subseteq$ in (10). Equality holds since $P_n+Q_n$ can have at most $2n$ vertices.
To infer (9) from (10), note that $u$ lies on the line connecting the origin and $v_{2n-1}=(2n-1,2n(2n-1))$.
\end{proof}
\begin{proof}[Proof of Theorem 30]
Inductively define a sequence of bivariate polynomials. The base case is
\[p_0=1\quad\text{and}\quad q_0=1.\]
The inductive step is
\[p_{n+1}=p_n(x^2,y^4)^2+x^Ny^Nq_n(x^2,y^4)^2\]
and
\[q_{n+1}=p_n(x^2,y^4)q_n(x^2,y^4)+x^{N(N-1)}y^{N^2(N-1)}\]
where $N=2^{n+1}$.
We claim that for every $n\ge0$,
\begin{equation}
\Newt(p_n)=2^nP_{2^n}\quad\text{and}\quad\Newt(q_n)=2^nQ_{2^n}.\tag{11}
\end{equation}
For $n=0$, this follows from $\Newt(p_0)=\Newt(q_0)=\{(0,0)\}=P_1=Q_1$. The induction step uses Lemma 11 and Observation 52. Assume that (11) holds for a given $n\ge0$. Then
\[\Newt(p_n(x^2,y^4))=2^n(2,4)P_{2^n}\quad\text{and}\quad\Newt(q_n(x^2,y^4))=2^n(2,4)Q_{2^n}.\]
Using (8),
\begin{align*}
\Newt(p_{n+1})&=2\cdot2^n(2,4)P_{2^n}\sqcup((N,N)+2\cdot2^n(2,4)Q_{2^n}))\\
&=2^{n+1}((2,4)P_{2^n}\sqcup((1,1)+(2,4)Q_{2^n})))\\
&=2^{n+1}P_{2^{n+1}}.
\end{align*}
Similarly, part (9) gives
\begin{align*}
\Newt(q_{n+1})&=2^n(2,4)(P_{2^n}+Q_{2^n})\sqcup\{(N(N-1),N^2(N-1))\}\\
&=2^{n+1}\big((1,2)(P_{2^n}+Q_{2^n})\sqcup\{(N-1,N(N-1))\}\big)\\
&=2^{n+1}Q_{2^{n+1}}.
\end{align*}
This proves (11).
To compute $p_n,q_n$, first construct a circuit of size $O(n)$ that simultaneously computes $x^M,x^{M(M-1)},y^M,y^{M^2(M-1)}$ for every $M=2^m$ with $m\le n$. Now, construct a circuit for $p_n$ and $q_n$ inductively. Given a circuit for $p_n$ and $q_n$, we can construct a new one computing $p_{n+1},q_{n+1}$ by introducing a constant number of extra gates.
\end{proof}
